\clearpage
\subsection*{C15. Reconstruction, geometric states and branch criteria}

\paragraph{From panorama endpoints to geometry.}
The panorama has width $W$ and height $H_{\mathrm{img}}$, with pixel coordinates
measured from the upper left. Each upper--lower pair first receives the
midpoint of its shorter horizontal arc, denoted $u$; its original vertical
coordinates $v_t$ and $v_b$ are retained. In radians, let
$\theta=2\pi(u/W-1/2)$ and
$\beta_j=\pi(v_j/H_{\mathrm{img}}-1/2)$ for $j\in\{t,b\}$.
The camera is the origin, the vertical axis points upward, and camera height
is the length unit. For $0<\beta_b<\pi/2$, the horizontal radius is
$\rho=1/\tan\beta_b$. The floor point is
$(\rho\sin\theta,-1,-\rho\cos\theta)$ and the corresponding upper point is
$(\rho\sin\theta,-\rho\tan\beta_t,-\rho\cos\theta)$.
Thus the corner's wall height is $h=1-\rho\tan\beta_t$.
The stored endpoint representation uses $W=1024$ and $H_{\mathrm{img}}=512$.
Measurements in metres, such as the depth example in the main paper, apply
the stated camera-height multiplier to these normalized coordinates.

\paragraph{Geometric validity and output states.}
The scoring geometry is an implicitly closed paired wall ring. Checks require
finite coordinates, equal upper and lower ring sizes, at least three pairs,
vertical pairing, the common floor plane, strictly positive wall heights,
and a simple nondegenerate footprint. Zero-length edges, self-intersections
and unresolved numerical degeneracies retain their failure states.
The spherical-boundary measurements additionally require defined finite
projections and nondegenerate shorter arcs. A valid floor can therefore
exist without a complete scorable wall layout. The scorer measures internal
orthogonal regularity and local height variation on this geometry; it does
not fit them away. Camera position and visibility-kernel diagnostics are
recorded separately, as in C14.

Scale-dependent numerical guards handle degenerate geometry and undefined projections; their implementation is specified in the accompanying source map.

\begingroup\small
\begin{longtable}{@{}>{\raggedright\arraybackslash}p{.24\linewidth}>{\raggedright\arraybackslash}p{.14\linewidth}>{\raggedright\arraybackslash}p{.56\linewidth}@{}}
\toprule
Branch & Minimum nodes & Additional condition and output\\
\midrule\endfirsthead
\toprule Branch & Minimum nodes & Additional condition and output\\\midrule\endhead
Geometric Q & 3 paired corners & Both assessed and reference layouts pass the geometric and metric-domain checks above. Valid zero floor overlap gives Q=0; unavailable geometry retains an unscorable status.\\[3pt]
Automatic consensus order & 3 selected identities & A projected source ring covers all selected identities with defined positions. Its connected geometry then undergoes the Q checks; a candidate order alone does not establish scorability.\\[3pt]
Fixed-coordinate assisted order & Same selected set & The supplied order preserves identities, support and coordinates. The resulting ring undergoes the same validity checks.\\[3pt]
Natural-size complete-output rates & 4 selected identities & A valid footprint gives a complete floor; paired wall geometry with recorded status \texttt{ok} gives the complete-evaluable geometry flag. Q additionally requires successful metric calculation.\\[3pt]
Source-history stopping & 4 nodes in a complete output & Starting at the requested prefix, add members until the recorded complete-geometry criterion is satisfied. Target Q does not choose the stopping point.\\
\bottomrule
\end{longtable}
\endgroup

These minimums belong to their respective procedures. The general geometry
and source-ring construction admit three nodes; the natural-size completion
flags and stopping delivery rule impose four.

\paragraph{Reference self-occlusion.}
For a valid reference footprint containing the camera strictly inside,
$H$ (\texttt{ref\_H}) counts floor vertices whose open sightline from the
camera crosses the interior of a non-incident footprint edge before reaching
the vertex. Tangential vertex contacts and collinear contacts are recorded
separately from crossings. Invalid footprints or a camera outside the
interior leave this feature unavailable. This is concealment by the room
boundary in the floor model. The same count enters the opportunity features
in C6 and the matched cross-view comparisons.

The accompanying source map in
\texttt{supplement\_sources/implementation\_reference/} identifies the
reconstruction, validity, scoring, correspondence and visibility functions
used to check these definitions.
